%================================ ANNEXES ===================================
% Chapitre numéroté en mode annexe (\appendix dans le fichier principal) :
% l'entrée « A Annexes » est ajoutée automatiquement à la table des matières,
% aucun \addcontentsline manuel n'est donc nécessaire (il créerait un doublon).
\chapter{Annexes}

\section{Commandes pour lancer le projet}

\subsection*{Backend (FastAPI)}
\begin{lstlisting}[language=bash]
cd backend
python -m venv .venv
# Windows PowerShell :
.\.venv\Scripts\Activate.ps1
# macOS / Linux :
# source .venv/bin/activate
pip install -r requirements.txt
uvicorn app.main:app --reload
# API      : http://127.0.0.1:8000
# Swagger  : http://127.0.0.1:8000/docs
# Santé    : curl http://127.0.0.1:8000/api/health
\end{lstlisting}

\subsection*{Frontend (React / Vite)}
\begin{lstlisting}[language=bash]
cd frontend
npm install
npm run dev
# Application : http://localhost:5173
npm run build   # build de production (dossier dist/)
npm run lint    # analyse statique ESLint
\end{lstlisting}

\section{Récapitulatif des endpoints de l'API}

L'API expose 24~endpoints sous le préfixe \texttt{/api} (auxquels s'ajoute un
point d'accueil \texttt{/} hors préfixe). Le tableau~\ref{tab:endpoints} en
présente la liste complète.

\begingroup\small
\begin{longtable}{@{}p{1.8cm}p{8.2cm}p{3.4cm}@{}}
\caption{Liste des endpoints de l'API backend}
\label{tab:endpoints}\\
\toprule
\textbf{Méthode} & \textbf{URL} & \textbf{Ressource} \\
\midrule
\endfirsthead
\toprule
\textbf{Méthode} & \textbf{URL} & \textbf{Ressource} \\
\midrule
\endhead
\bottomrule
\endfoot
GET & \path{/api/health} & Santé \\
GET & \path{/api/statistics/dashboard} & Tableau de bord \\
GET & \path{/api/campaigns} & Campagnes \\
POST & \path{/api/campaigns/drafts} & Brouillons \\
GET & \path{/api/campaigns/drafts} & Brouillons \\
GET & \path{/api/campaigns/drafts/{draft_id}} & Brouillons \\
GET & \path{/api/campaign-creation/options} & Options de l'assistant \\
GET & \path{/api/stores} & Magasins \\
POST & \path{/api/stores/import/preview} & Import de magasins \\
POST & \path{/api/dco/creatives} & DCO \\
GET & \path{/api/dco/creatives} & DCO \\
GET & \path{/api/advertisers} & Annonceurs \\
GET & \path{/api/advertisers/{advertiser_id}} & Annonceurs \\
PATCH & \path{/api/advertisers/{advertiser_id}} & Annonceurs \\
GET & \path{/api/advertisers/{advertiser_id}/settings} & Paramètres annonceur \\
PATCH & \path{/api/advertisers/{advertiser_id}/settings} & Paramètres annonceur \\
GET & \path{/api/account-settings} & Paramètres globaux \\
PUT & \path{/api/account-settings} & Paramètres globaux \\
PATCH & \path{/api/account-settings} & Paramètres globaux \\
GET & \path{/api/users} & Utilisateurs \\
POST & \path{/api/users} & Utilisateurs \\
GET & \path{/api/users/{user_id}} & Utilisateurs \\
PATCH & \path{/api/users/{user_id}} & Utilisateurs \\
PATCH & \path{/api/users/{user_id}/status} & Utilisateurs \\
\end{longtable}
\endgroup

\section{Checklist de démonstration (extrait)}

\textbf{Avant la démonstration :}
\begin{itemize}
  \item redémarrer le backend pour repartir sur des données propres ;
  \item vérifier \texttt{GET /api/health} ;
  \item lancer le frontend et ouvrir l'application dans un onglet rafraîchi
        (Ctrl+F5) ou en navigation privée ;
  \item ouvrir Swagger dans un second onglet ;
  \item préparer le fichier d'exemple \path{docs/samples/stores-valid.csv}.
\end{itemize}

\textbf{Incidents courants et solutions :}
\begin{itemize}
  \item \textbf{backend non lancé} : relancer
        {\NoAutoSpacing\texttt{uvicorn app.main:app --reload}} ;
  \item \textbf{port occupé} : identifier et arrêter le processus concerné,
        puis relancer le serveur ;
  \item \textbf{page blanche ou cache obsolète} : recharger en forçant le
        cache (Ctrl+F5) ou vider \path{node_modules/.vite} avant de relancer ;
  \item \textbf{export CSV mal affiché} : ouvrir le fichier via l'import de
        données d'Excel si le double-clic pose problème ;
  \item \textbf{erreur API} : reproduire l'appel dans Swagger pour lire le
        message d'erreur exact.
\end{itemize}

Le détail complet figure dans le fichier
\path{docs/checklist-demonstration-finale.md} du dépôt.

\section{Flux de travail Git}

\begin{lstlisting}[language=bash]
# Créer une branche par sujet
git checkout -b feat/mon-module

# Valider les changements (commit)
git add .
git commit -m "feat: description du module"

# Pousser la branche et ouvrir une pull request sur GitHub
git push -u origin feat/mon-module
# -> ouverture d'une pull request, revue, puis fusion (merge)
\end{lstlisting}

\section{Liste des captures d'écran}

\begin{enumerate}
  \item Interface de connexion (figure~\ref{fig:login})
  \item Tableau de bord principal (figure~\ref{fig:dashboard})
  \item Liste des campagnes (figure~\ref{fig:campagnes})
  \item Assistant de création de campagne (figure~\ref{fig:creation})
  \item Import des magasins (figure~\ref{fig:magasins})
  \item Carte de géociblage (figure~\ref{fig:geofencing})
  \item Module DCO (figure~\ref{fig:dco})
  \item Reporting des performances (figure~\ref{fig:reporting})
  \item Gestion du compte (figure~\ref{fig:compte})
  \item Documentation Swagger de l'API (figure~\ref{fig:swagger})
  \item Test de l'endpoint \texttt{/api/health} (figure~\ref{fig:swaggerhealth})
  \item Dépôt GitHub (figure~\ref{fig:github})
\end{enumerate}
